%%%%%%%%%%%%%%%%%%%%%%%%%%%%%%%%%%%%%%%%%%%%%%%%%%%%%%%%%%%%
% RELATED WORKS
%%%%%%%%%%%%%%%%%%%%%%%%%%%%%%%%%%%%%%%%%%%%%%%%%%%%%%%%%%%%
\section{Related Works}
Our work builds on the extensive literature of three areas: forward PDE solvers, inverse PDE solvers, and diffusion models.
Please see relevant surveys for more information \cite{evans2022partial, omori2007overview, po2023state,strauss2007partial, yang2023diffusion}.

\vspace{-1.5em}\noindent
\paragraph{Forward PDE Solvers.} 
PDE solvers take the specification of a physics system and predict its state in unseen space and time by solving an equation involving partial derivatives.
Since Most PDEs are very challenging to solve analytically, people resolve to numerical techniques, such as Finite Element Method \cite{quarteroni2008numerical,solin2005partial} and Boundary Element Method~\cite{aliabadi2020boundary,idelsohn2003meshless}.
While these techniques show strong performance and versatility in some problems, they can be computationally expensive or difficult to set up for complex physics systems.
Recently, advancements in deep-learning methods have inspired a new set of PDE solvers.
Raissi et al. \cite{raissi2017physics, raissi2019physics} introduce Physics-Informed Neural Networks (PINNs), which optimize a neural network using PDE constraints as self-supervised losses to output the PDE solutions. 
PINNs have been extended to solving specific fluid \cite{cai2021physicsfluid, mao2020physics}, Reynolds-averaged Navier–Stokes equations \cite{eivazi2022physics}, heat equations \cite{cai2021physicsheat}, and dynamic power systems \cite{misyris2020physics}.
While PINNs can tackle a wide range of complex PDE problems, they are difficult to scale due to the need for network optimization.
An alternative approach, neural operators~\cite{li2020fourier,lu2021learning}, directly learn the mapping from PDE parameters (\eg initial and boundary condition) to the solution function.
Once trained, this method avoids expensive network optimization and can instantly output the solution result.
This idea has been extended to solve PDE in 3D~\cite{li2024geometry,serrano2024operator} , multiphase flow \cite{wen2022u}, seismic wave \cite{lehmann2023fourier,li2023solving}, 3D turbulence \cite{li2022fourier, peng2023linear}, and spherical dynamics \cite{bonev2023spherical}. 
People have also explored using neural networks as part of the PDE solver, such as compressing the physics state~\cite{chen2022crom,li2023neural,muller2021real,nam2024solving}.
These solvers usually assume known PDE parameters, and applying them to solve the inverse problem can be challenging.

\vspace{-1.5em}\noindent
\paragraph{PDE inverse problem.}
The inverse problem refers to finding the coefficients of a PDE that can induce certain observations, mapping from the solution of a PDE solver to its input parameters.
People have tried to extend traditional numerical methods to this inverse problem~\cite{cho2017first,fox2001statistical,harrach2021introduction, kumar2023accelerating,van2015penalty}, but these extensions are non-trivial to implement efficiently.
There are similar attempts to inverse deep-learning PDE solvers.
For example, one can inverse PINNs by optimizing the network parameters such that their outputs satisfy both the observed data and the governing equations.
iFNO \cite{long2024invertible} and NIO \cite{molinaro2023neural} tries to extend FNO~\cite{li2020fourier}. 
Other methods \cite{de2022deep,pakravan2021solving} directly learn the operator functions for the inverse problem. 
PINO \cite{li2021physics} further combines neural operators with physics constraints to improve the performance of both forward and inverse problems.
These methods assume full observations are available.
To address the inverse problem with partial observations, people have tried to leverage generative priors with Graph neural networks~\cite{iakovlev2020learning, zhao2022learning}.
These works have not demonstrated the ability to solve high-resolution PDEs, possibly limited by the power of generative prior. 
We want to leverage the state-of-the-art generative model, diffusion models, to develop a better inverse PDE solver.

\vspace{-1.5em}\noindent
\paragraph{Diffusion models.}
Diffusion models have shown great promise in learning the prior with higher resolutions by progressively estimating and removing noise. 
Models like DDIM \cite{song2020denoising}, DDPM \cite{ho2020denoising}, and EDM \cite{Karras2022edm} offer expressive generative capabilities but face challenges when sampling with specific constraints. 
Guided diffusion models \cite{chung2022diffusion,chung2022improving,dhariwal2021diffusion,xu2023inversion} enhance generation processes with constraints such as image inpainting, providing more stable and accurate solutions. 
Prior works on diffusion models for PDEs highlight the potential of diffusion approaches by generating PDE datasets such as 3D turbulence \cite{jacobsen2023cocogen,lienen2023zero} and Navier-Stokes equations \cite{yang2023denoising} with diffusion models. Diffusion models can also be used to model frequency spectrum and denoise the solution space \cite{lippe2024pde}, and conditional diffusion models are applied to solve 2D flows with sparse observation \cite{shu2023physics}. However, the application of diffusion models to solve inverse problems under partial observation remains underexplored. In this work, we aim to take the initial steps towards addressing this gap.
